
\section{Inverse Reinforcement Learning\label{sec:irl}}

The IRL literature was initiated by \citet{ng:2000} who identified two motivations for inferring the underlying
rewards that are presumed to generate observed behavior: 1) ``the potential use of reinforcement learning and
related methods as potential models for human and animal learning'' and 2) to provide a reward function
that RL can use to construct an ``intelligent agent that can behave successfully in a particular domain.'' Though
they note that imitation and apprenticeship learning (IL and AL) can be used to learn a policy $\pi$, ``the reward
function often provides a much more parsimonius description of behaviour'' so that IRL can constitute an effective
form of AL. IRL has become an important method in many areas of robotics including development of autonomous 
drones and vehicles (AV) where there is insufficient data on decisions to human drivers under multitudes of different
road conditions to develop effective AVs.  \citet{av:2025} note that IRL has ``found great success in predicting 
trajectories through inferring cost or reward functions from expert demonstrations, 
and then using these functions to guide the behavior of self-driving vehicles in unseen driving environments.''

This section provides a selective review of the IRL literature focusing on methods that overlap
most closely with the structural DDC literature reviewed in section 3. We refer readers to surveys 
by \citet{arora:2021}, \citet{gleave:2022}, and \citet{adams:2022} for a broader and more in depth view of IRL.
The latter survey
provides further examples of how IRL uses  ``an expert to demonstrate the successful task and then a highly parameterized reward function could be transferred to the robot for optimization'' and how IRL can generate more effective policies than AL, particularly
for new, counterfactual environments where
``IRL has the added advantage that the reward function can be transferred to new environments with different dynamics.''.

There are three broad classes of IRL algorithms: 1) max-margin, 2) Bayesian methods, and 3) maximum entropy methods.
\citet{ng:2000} introduced three different types of max-margin estimators, which
search for a reward function for which payoffs under the observed policy are 
higher than any other policy (i.e. behavior), an approach that shares similarities with the ``moment inequality'' simulation
estimator  proposed by \citet{bbl:2007} discussed in section 3. \citet{ng:2000} recognized the identification problem discussed
in the previous section, namely there may be many reward functions that satisfy the max-margin inequalities (payoff inequalities)
including the trivial solution $r=0$. 

Bayesian IRL (BIRL) was introduced by \citet{birl:2007}. They assumed $\pi$ is a mixed strategy with the softmax
form (\ref{eq:ccp}) that depends on $Q$ functions which are implicit functions or $r$ per the Bellman equations (\ref{eq:q1}) and
(\ref{eq:q2}) and interpret the $\sigma$ parameter in the softmax function for $\pi$ as
 controlling for the ``degree of confidence we have in the [DM's] ability to
choose actions with high value.'' This allowed them to formulate a likelihood for the sequence of observed actions by
the DM that depends on $r$. Using  Bayes' Rule, they computed a posterior distribution for $r$, treating it as the unknown parameter
with a prior distribution. Since it is infeasible to directly compute the
posterior,  BIRL uses MCMC to generate draws from the posterior
distribution, with policy iteration used to solve for $Q$ for each draw of $r$. This makes BIRL method computationally intensive
for large state spaces similar to NFXP. Overall BIRL is similar to the Bayesian MCMC methods discussed in section 3, except that those
methods use stochastic updates of $Q$ similar to those produced by Q-learning to gradually compute $Q$ over the course
of the MCMC simulations rather than use policy iteration to find $Q$ for each MCMC draw of $r$.

A seminal contribution to the IRL literature is
\citet{ziebart:2008}, who introduced {\it maximum entropy IRL\/} as a means of justifying the use of the softmax/logit
probability for  $\pi(a|s)$ as the probability with maximum entropy subject to certain constraints that make
expectations of certain functions known as {\it features\/} equal the sample averages of these features.\footnote{\footnotesize
\citet{ziebart:2010} extended these ideas to MDPs and called it {\it Maximum causal entropy RL.\/}}
They proposed to estimate the parameters of a reward function that is a linear
combination of a set of known features using a partial likelihood for observed $(a,s)$ pairs
that is identical to the partial likelihood (\ref{eq:ccplf}) used in the DDC literature.\footnote{\footnotesize
The full likelihood
(\ref{eq:lf}) is not frequently used as the estimation criterion in IRL
 because of the model-free or non-parametric approach that bypasses
the need to explicity specify or estimate the transition probability $p(s'|s,a)$.} 
They applied this method to model the route choices of 25 Yellow cab drivers in Pittsburgh with a reward
function that involves four features: road type, speed, lanes and transitions. They showed that the estimated
model could correctly predict nearly 80\% of the cab paths in a holdout or evaluation sample and that
their new method was a significant improvement over other IRL methods such as max-margin.

These ideas evolved into the already cited work of \citet{softqlearning:2017} who characterized $\pi$ as the policy that maximizes the
expected discounted entropy-regularized payoff given in equation (\ref{eq:vdefentropy}). 
Various criteria are used to estimate $r$ such as ``occupancy matching'' (\citet{hoermon:2016}, \citet{clare:2023}). For our purposes
we focus on IRL where the likelihood function is the estimation criterion, such as \citet{ermon:2015}, \citet{zeng:2023}, \citet{zeng:2024},
\citet{rp:2022} and \citet{kang:2025}.\footnote{\footnotesize \citet{zeng:2023} and \citet{zeng:2024} maximize a discounted log-likelihood function.} 
This results in a nearly complete equivalence between the DDC and IRL literatures, with
identical softmax/logit formulas for the policy $\pi(a|s)$ (\ref{eq:ccp}) 
that depend on the choice-specific or ``soft $Q$'' values that satisfy the ``soft Bellman equation''
in equation (\ref{eq:smoothQ}). Overall, DDC and maximum entropy IRL result in
mathematically identical optimization problems for recovering $r$ from data on states and actions of a sample of DMs.
The main difference is in the algorithms used to solve the DM's DP problem.

In DDC the $Q$ functions are computed  by non-stochastic algorithms such as SA or PI, or in large state spaces
by parametrizing $Q$ or the value function $V$ and finding parameters $\phi$ that minimize the squared ``Bellman
residuals'' as in the SEES estimator, (\ref{eq:penalizedllf}). All of these methods can be described as {\it model-based\/}
because they require knowledge of the transition probability $p(s'|s,a)$, or at least $p$ must be estimable from
data. \citet{adams:2022} notes that most IRL methods are model-based. However, we will focus on 
model-free IRL that does not require specification of a function form for $p$ or
even require that it be explicitly estimated non-parametrically. Instead, 
new offline IRLs inspired model-free methods such as Q-learning are able to approximate the $Q$ functions 
using only realizations from $p$ without knowing $p$ itself.  
But how is it possible to generate realizations $s'$ from $p(s'|s,a)$ without knowing $p$ itself?

Model-free IRL solves this problem in a clever way: it uses already realized transitions in demonstration datasets to provide the
``simulator'' needed by RL algorithms  to train the Q and values functions.
The insight is that the data on observed states and decisions provides the  ``environment'' and the state transitions needed
by RL to find  the value and Q functions needed to recover $r$, reminiscent of the idea
of {\it experience replay\/} in the deep Q learning literature, \citet{mnih:2015}.\footnote{\footnotesize Also related is the
distinction between online and offline RL. IRL generally uses offline RL since the estimation of rewards is based
on already collected datasets rather than being done in real time.  See the survey of offline RL by
 \citet{offlinerl:2020} who note ``Offline reinforcement learning algorithms hold tremendous promise for making it possible to turn large datasets into powerful decision making engines.''.  \citet{clare:2023} note that offline IRL also focuses on 
``learning from a previously collected dataset without online interaction with the environment'' 
and ``holds tremendous promise for safety-sensitive applications where manually identifying an appropriate reward 
is difficult but historical datasets of human demonstrations are readily available (e.g., in healthcare, autonomous driving, robotics, etc.).''}
To get a better understanding of these methods, we will provide a more in-depth discussion of two recent contributions the IRL literature
that can recover $r$ from data in a model-free manner. One can be regarded as an IRL version of the CCP estimator and
the other is similar to the SEES estimator (\ref{eq:penalizedllf}).  Both of these
methods use parametric function approximation methods (expressing $V$ and $Q$ as outputs of neural networks or
linear combinations of a set of basis functions) that enables them to handle problems with large 
state spaces.


\citet{ae:2022} (AE) introduce a model-free two step CCP estimator similar
the one  described in the previous section but using RL-inspired methods to
avoid estimating or making parametric assumptions about $p(s'|s,a)$. They do assume 
that the reward function is a linear combination of known ``features'', i.e. $r_\theta(s,a)=\sum_{k=1}^K \theta_k r_k(s,a)$.
Their estimator is inspired by the temporal difference (TD) method of RL,
and they apply it to approximate the functions $R$ and $Q_\epsilon$ 
(defined in equations (\ref{eq:Qeps}) and (\ref{eq:hequation}), respectively) that
enter as covariates in the CCP $\pi$ in equation (\ref{eq:ccpestimator}), using
transitions in the data to estimate these functions 
without having to explicitly estimate $p(s'|s,a)$. To  accommodte 
empirical applications with large state spaces they propose parameterizing $R$ and $Q_\epsilon$ as linear
combinations of a set of known basis functions that depend on $(s,a)$. Let $\nu(s,a)=\{\nu_1(s,a),\ldots,\nu_J(s,a)\}$
denote $J$ basis functions that are used to approximate $R$ nad $Q_\varepsilon$ by a regression that uses the
basis functions evaluated at values of $(s,a)$ in the data set as regressors.  We will describe AE's approach
to approximating $R$, since the same procedure is used to approximate $Q_\varepsilon$ and is omitted
for brevity. Let  $R_\phi$ denote an approximation to $R$ based
on a $J \times 1$ parameter vector $\phi$ of the form $R_\phi(s,a)=\nu(s,a)\phi=\sum_{j=1}^J \theta_j\nu_j(s,a)$. AE propose to 
to choose $\hat\phi$ to minimize the mean-squared TD error given by
\begin{equation}
     \hat\phi =\argmin_\phi E\left\{ [\vec{r}(s,a)+\beta R_\phi(s',a')-R_\phi(s,a)]^2\right\},
\label{eq:hatphi}
\end{equation}
where $E$ denotes the expectation operator over $(s,a,s',a')$.
Using results from \citet{TsitsiklisVanRoy:1997} on the convergence of TD learning with linear function approximation, 
AE show that $\hat\phi$ can be computed using the sample expectation operator $E_N$ in (\ref{eq:hatphi})
using the formula
\begin{equation}
     \hat\phi = \left[E_N\{\nu(s,a)[\nu(s,a)-\beta\nu(s',a')]^{{\scriptscriptstyle T}}\}\right]^{-1} E_N\{\nu(s,a)\vec{r}(s,a)\},
\label{eq:QR_regression}
\end{equation}
and the implied estimate $\hat R(s,a)=\phi(s,a)\hat\phi$ of the features or 
``covariates'' for the $\theta$ in (\ref{eq:ccpestimator}). We also have to approximate the 
second ``correction term''  $\hat Q_\varepsilon(s,a)$, which can be estimated via  regression
similar to (\ref{eq:QR_regression}) but it requires a first stage non-parametric estimate of $\pi(a|s)$ to construct
estimates of the expected choice-specific unobserved shocks $E\{\varepsilon(a)|s,a\}=-\log(\pi(a|s))$.\footnote{\footnotesize 
AE normalize  $\sigma$ to 1 since it cannot be separately identified from the $\theta$ parameters.}
Using the estimated $\hat R(s,a)$ and $\hat Q_{\varepsilon}(s,a)$ values, 
 they can compute a second stage pseudo-maximum likelihood (PMLE) estimator of $\theta$ by maximizing 
the partial likelihood (\ref{eq:ccplf})  using the formula for $\pi_\theta(a|s)$ in (\ref{eq:ccpestimator}). Consistency
of this first stage estimator requires specification of a sieve with $J$ basis functions $\{\nu_1,\ldots,\nu_J\}$ that increases with
sample size $N$ at an appropriate rate and with sufficient linear independence
to allow any function $R$ to be perfectly approximated by its projection $\Pi_\nu R$ in the limit as $J \to \infty$. 
The PMLE is affected by ``estimation noise'' in the $\hat R$ and $\hat Q_{\varepsilon}$ functions
that can introduce bias and noise. To deal with this, AE propose a third estimation stage that uses the PMLE
to create a modified  score (i.e. gradient) of the  partial log-likelihood (\ref{eq:ccplf}) to create 
a robust (Neyman orthogonal) version of the score function
which they used to their preferred estimator $\hat\theta$  as the value that sets the modified score to zero following the
approach of \citet{lrspe:2022}. 

AE proposed a different  estimator they call AVI (for Approximate Value Iteration) that allows for nonlinear parametrizations
of the $R$ and $Q_\varepsilon$ functions which are computed using an iterative procedure for updating the parameter $\hat\phi_k$ is the estimated
parameters determining $\hat R_k(s,a)=R_{\hat\phi_k}(s,a)$ as follows
\begin{equation}
    \hat\phi_{k+1} = \argmin_\phi E_N\left\{\left[ \vec{r}(s,a)+\beta R_{\hat\phi_k}(s',a')-R_{\phi}(s,a)\right]^2\right\}.
\label{eq:avi}
\end{equation}
These  iterations continue  until the sequence $\{\hat\phi_k\}$ converges and the resulting estimator of $R$ is $\hat R=R_{\hat\phi_k}$ where
$\hat\phi_k$ is the last value of the iterative procedure (\ref{eq:avi}), and similarly for  $\hat Q_\varepsilon$.\footnote{\footnotesize 
It is possible to extend their approach to reward functions $r_\theta(s,a)$ where $\theta$ does not enter as a linear combination
of a set of known features $\vec{r}(s,a)$. However in that case, we can no longer express $Q_r(s,a)=R(s,a)*\theta$, and in general
$\theta$ will enter $Q_r$ in a non-linear fashion. The TD estimator can still be adapted to work in this situation except
that now there must be a nested TD subroutine that recomputes $\hat Q_r$ as a function of $\theta$ each time $\theta$ is updated
in an outer optimization algorithm that searches over $\theta$ to maximize the partial likelihood function (\ref{eq:ccplf}).}
AE conclude that ``A range of Monte Carlo simulations using a dynamic firm entry problem, a dynamic firm entry game and two versions of the famous Rust (1987) engine replacement problem show that the proposed algorithms work well in practice.''\footnote{\footnotesize
See \citet{khwaja:2025} for an alternative approach to AE's TD methods that they call RLTD-CCS that combines the TD algorithm 
from RL with the Conditional Choice Simulator (CCS) estimator of \citet{hmss:1994} that relies on forward simulations
of paths to produce a Monte Carlo estimate the $Q_\theta(s,a)$  and hence is a model-based approach to estimation.}

\citet{kang:2025} developed an estimator they call GLADIUS (Gradient-based Learning with Ascent-Descent for Inverse Utility learning from Samples)
that is 1) model-free, 2) does not require a first stage non-parametric estimation of $\pi(a|s)$,
and 3) does not require a parametric specification for rewards. Instead, GLADIUS is an iterative 
algorithm that uses gradient ascent/descent to maximize a penalized likelihood function similar to the SEES estimator
(\ref{eq:penalizedllf}). However instead of parametrizing $r$ as a function of ``structural parameters''
$\theta$  and using using auxilliary parameters $\phi$ to approximate
the $Q$ functions, GLADIUS relies on 
flexible parametrization of the $Q$ function as $Q_{\theta_1}(s,a)$ and the conditional expectation of the value function $E\{V(s')|s,a\}$ as
$EV_{\theta_2}(s,a)$ for parameters $\theta_1$ and $\theta_2$ of flexible approximations
to $Q$ and $EV$ such as  DNNs. Once GLADIUS converges, returning parameters $\hat\theta_1$ and
$\hat\theta_2$, rewards are estimated as the function
\begin{equation}
             r(s,a) = Q_{\hat\theta_1}(s,a)-\beta EV_{\hat\theta_2}(s,a),
\end{equation}
i.e. they are  ``backed out'' from the fixed point equation for smoothed $Q$ function, (\ref{eq:smoothQ}).

Both GLADIUS and SEES use an objective function that maximizes a penalized likelihood where the penalty term 
is a ``mean squared Bellman error'' for the $Q$ functions. That is, if we express $Q$ as a fixed point $Q=\Lambda_\sigma(Q)$,
$Q$ minimizes the quantity $\rho_{BE}(Q)=E\{[Q(s,a)-\Lambda_\sigma(Q)(s,a)]^2\}$ over a joint distribution of $(s,a)$ pairs.
SEES is a model-based approach that uses an assumed parametric functional form
for rewards, $r_\theta(s,a)$ and  knowledge (or estimated version thereof) of the transition probability'
$p(s'|s,a)$ to compute the Bellman operator $\Lambda_{\theta,p,\sigma}$, where we subscript
$\Lambda_\sigma$ to emphasize its dependence on $\theta$ as well as $p$. Thus, SEES calculates
the mean squared Bellman error $\rho_{BE}(Q)$ numerically using deterministic methods such as numerical quadrature
to compute the conditional expectations of the $Q$ functions. 

GLADIUS is a model-free approach inspired
by Q-learning that uses transitions $(s',s,a)$ in the data to construct an 
estimate of the mean squared Bellman error $\rho_{BE}(Q)$. That is, GLADIUS seeks to approximate the quantity
\begin{equation}
\Lambda_\sigma(Q)(s,a)=r(s,a)+\beta \sum_{s'} V(s')p(s'|a,d)
\label{eq:Lambda}
\end{equation}
by averaging realizations of the form 
\begin{equation}
   \tilde\Lambda_\sigma(Q)(s',s,a) = r(s,a)+\beta V(s'),
\label{eq:Lambdadraw}
\end{equation}
using observed $(s,a,s')$ transitions in the sample. It follows that
$E\{\tilde\Lambda_\sigma(s',s,a)|s,a\}=\Lambda_\sigma(Q)(s,a)$ so 
we can use transitions $(s',s,a)$ in the data to calculate
 $\Lambda_\sigma(Q)(s',s,a)$ values, and then use these as dependent variables
in a regression to estimate its conditional expectation  $\Lambda_\sigma(s,a)$.
This suggests a strategy of  estimating $Q$ by searching for $\theta_1$ to minimize
the mean squared temporal difference error $\rho_{TD}(Q)$
given by
\begin{equation}
        \rho_{TD}(Q)=E\{[\Lambda_\sigma(Q)(s',s,a)-Q(s,a)]^2\},
\end{equation}
where, using the terminology of the RL literature,  $\tilde\Lambda_\sigma(Q)(s',s,a)-Q(s,a)$ is the TD error.

However \citet{kang:2025} note an important  complication: the true fixed point $Q$ will minimize
the mean squared Bellman error $\rho_{BE}(Q)$ but it will not necessarily minimize the mean squared
TD error. To see, this, the note the following identity
\begin{equation}
   \rho_{TD}(Q)=\rho_{BE}(Q)+\beta^2 E\{[V(s')-EV(s,a)]^2\},
\label{eq:tdbe}
\end{equation}
where the second term reflects expected ``heteroscedasicity'' in the TD error, i.e. it is the average of the
conditional variances given by $H(V)(s,a)=E\{[V(s')-EV(s,a)]^2|s,a\}$. Since $V$ is a function of $Q$
from (\ref{eq:smoothV}), we write $V$ as $V(Q)$ and the heteroscedasticity term as $H(Q)$.  It 
it follows from (\ref{eq:tdbe}) that the $Q$ that minimizes $\rho_{TD}(Q)$ will not necessarily be the same as the $Q$ that minimizes $\rho_{BE}(Q)$.
To deal with this issue, they express the mean squared Bellman
error as  $\rho_{BE}(Q)=\rho_{TD}(Q)-H(Q)$. Though $H(Q)$ is also an unknown quantity
Kang {\it et. al.\/}  note that it can be approximated as the solution to
\begin{equation}
             H(Q)(s,a) =\min_{v \in R} \sum_{s'}[V(Q)(s')-v]^2 p(s'|s,a).
\label{eq:condvar}
\end{equation}
Of course the $v$ that minimizes (\ref{eq:condvar}) is  $EV(s,a)$.  \citet{kang:2025} parametrize $Q$ and $EV$  by DNNs as 
$Q_{\theta_1}$ and $EV_{\theta_2}$, respectively, to obtain an estimation criterion that is a penalized likelihood
function similar to SEES, except it is a max-min problem  given by
\begin{equation}
       (\hat\theta_1,\hat\theta_2) =\argmax_{\theta_1}\argmin_{\theta_2} \log(L(\theta_1))-\lambda\rho(\theta_1,\theta_2),
\label{eq:maxmin}
\end{equation}
where $\lambda$ is a penalization parameter, and $\rho(\theta_1,\theta_2)$ is the estimate of the mean squared Bellman error
given by
\begin{equation}
       \rho(\theta_1,\theta_2) =\frac{1}{N} \sum_i^N [\Lambda_\sigma(Q_{\theta_1}(s'_i,s_i,a_i)-Q_{\theta_1}(s_i,a_i)]^2 - \beta^2
 [V(Q_{\theta_1})(s'_i)-EV_{\theta_2}(s_i,a_i)]^2 
\label{eq:kangpenality}
\end{equation}
Computation of the penalty term can be further reduced to limiting the summation to a subset of states
$s_i$ where the corresponding action $a_i$ equals the {\it normalizing action.\/} The reason is that
\cite{kang:2025} impose the identifying assumption that for each $s$ there is an action $a_s$ called the
normalizing action for which $r(s,a_s)$ is known. From our discussion of identification in section 3.2, this
restriction on rewards implies that we can determine all $Q$ values by enforcing the fixed point
constraint $Q=\Lambda_\sigma(Q)$ only for the subset of pair $(s,a_s)$ which enables us to determine the
corresponding $Q$ at the normalizing action,  $Q(s,a_s)$ for all $s \in S$. Then given $Q(s,a_s)$ the remaining
$Q(s,a)$ values are determined from $\pi(a|s)$ using the Hotz-Miller inversion  formula (\ref{eq:ccpinverse}).

\citet{kang:2025} establish an important result on the global convergence of GLADIUS. First, under a ``realizability assumption''
that there exists $\theta_1^*$ such that $Q^*=Q_{\theta_1^*}^*$, where $Q^*$ is the true Q function
(as well as some other regularity conditions), they show that 
if $\hat Q_{T,N}$ is the estimate of $Q^*$ after $T$ iterations of the GLADIUS algorithm using a data set with $N$ observations
on agent's states and transitions $(s',s,a)$, then
\begin{equation}
     E_N\left\{[Q^*(s,a)-\hat Q_{T,N}(s,a)]^2\right\} \le O(1/T) + O(1/N).
\end{equation}
They note that ``To the best of our knowledge,  
no prior work has proposed an algorithm that guarantees global optimum convergence of the minimization problem that involves
the [mean squared Bellman residual term]''. They compare GLADIUS to several other DDC estimation methods using the
\citet{rust:1987} bus engine problem and a high dimensional version of this problem where they added between
2 and 100 additional irrelevant
state variables (i.e. variables that have no effect on cost of maintaining the bus or its mileage
transitions, and thus on the Q values and decision rule $\pi$) that take on 20 possible values each, thus 
massively increasing the size of the state space. They conclude that
``We find that, on average, our method performs quite well in recovering rewards in both low and high-dimensional settings. Further, it has better/on-par performance compared to other benchmark algorithms in this area (including algorithms that assume the parametric form of the reward function and knowledge of state transition probabilities) and is able to recover rewards even in settings where other algorithms are not viable.''\footnote{\footnotesize
\citet{kang:2025} also note GLADIUS corrects a problem in IRL methods that use ``occupancy matching'' as the estimation criterion
such as \citet{hoermon:2016}, \citet{iqlearn:2022}, and \citet{clare:2023}. They show that occupancy matching is not guaranteed
to consistency estimate the true $Q$ function, and ``This implies that $r$ cannot be inferred from $Q$ 
using the Bellman equation after deriving $Q$ using occupancy matching.''}

We have provided a very selective survey of the 
IRL literature, going into depth on two recent contributions to this literature
that are particularly closely related to and we think particular promising new methods 
for structural estimation of DDC models.  We have emphasized the distinction between model-based and model-free
estimation methods, but we wish to make clear that not all IRL methods are model-free. 
For example \citet{zeng:2024} use non-parametric methods to estimates $p$ to provide a ``generative world model'' 
in a model-based version of IRL and ``Through extensive experiments, we demonstrate that our algorithm outperforms existing benchmarks for offline IRL and Imitation Learning, especially on high-dimensional robotics control tasks.'' 
Thus, it is not obvious that model-free approaches are necessarily superior to model-based approaches,
though we agree that good models of $p$ should be used, since  ``inaccurate models of the world obtained from finite data with limited coverage could compound inaccuracy in estimated rewards.''  It is also clear that if we are interested in 
doing a counterfactual that involves changing $p$, model-based methods will
probably be necessary to estimate/construct this new $p$ and
then train the model to calculate the new optimal decision rule $\pi(a|s)$ corresponding to counterfactual $p$.

We also do not want to leave an impression that the two studies we have focused on are the first to use
function approximation methods such as DNNs to estimate models with large scale state spaces
in the IRL literature. There are also many seemingly different variants of
IRL that are actually closely related on closer inspection. For example, \citet{ziebart:2010}
introduced {\it maximum causal entropy\/} (MCE) IRL, which is nearly isomorphic to DDC models
under the AS/CI/EV assumptions discussed in  section 3.  As the survey by \citet{gleave:2022} notes:
``algorithms based on MCE IRL have scaled to high-dimensional environments. Maximum Entropy Deep IRL \citet{maxentdeepirl:2016} was one of the first extensions, and is able to learn rewards in gridworlds from pixel observations. More recently, Guided Cost Learning 
\citet{guidedcostlearning:2016} and Adversarial IRL \citet{airl:2018} have scaled to MuJoCo continuous control tasks.''  
The introduction also noted two large scale studies using IRL: the study by
\citet{barnes:2024} that uses a version of IRL called Receding Inverse Horizon Planning to infer preferences  over
route choice using data from Google Maps,  and  the study by \citet{zhao:2023} who used AIRL and showed 
that it outperforms competing methods in an application to inferring context-dependent preferences for route choice 
using data on taxi drivers in Shanghai. All this work indicates the promise of IRL methods for estimation
of large scale DDC problems.
 

